\documentclass[11pt]{article}

\usepackage[T1]{fontenc}
\usepackage[utf8]{inputenc}
\usepackage[margin=1in]{geometry}
\usepackage{amsmath,amssymb,amsthm,mathtools}
\usepackage{booktabs}
\usepackage{array}
\usepackage{longtable}
\usepackage{enumitem}
\usepackage{xcolor}
\usepackage{hyperref}
\usepackage[nameinlink,capitalise]{cleveref}

\hypersetup{
  colorlinks=true,
  linkcolor=blue!45!black,
  citecolor=blue!45!black,
  urlcolor=blue!45!black,
  pdftitle={Named Grid Covers under Row--Column Margins: Three Exactness Levels and Largest-Piece Localization},
  pdfauthor={Oleksiy Babanskyy}
}
\setlist{nosep}
\setlength{\emergencystretch}{2em}

\ifdefined\pdfinfoomitdate
  \pdfinfoomitdate=1
  \pdfsuppressptexinfo=-1
  \pdftrailerid{}
\fi

\newcommand{\ZZ}{\mathbb Z}
\newcommand{\NN}{\mathbb N}
\newcommand{\TT}{\mathcal T}
\newcommand{\FFib}{\mathcal F}
\newcommand{\PP}{\mathcal P}
\newcommand{\EE}{\mathcal E}
\newcommand{\GG}{\mathcal G}
\newcommand{\KK}{\mathcal K}
\newcommand{\MM}{\mathcal M}
\newcommand{\one}{\mathbf 1}
\newcommand{\supp}{\operatorname{supp}}
\newcommand{\rank}{\operatorname{rank}}
\newcommand{\dist}{\operatorname{dist}}
\newcommand{\im}{\operatorname{im}}
\newcolumntype{L}[1]{>{\raggedright\arraybackslash}p{#1}}

\theoremstyle{plain}
\newtheorem{theorem}{Theorem}[section]
\newtheorem{proposition}{Proposition}[section]
\newtheorem{lemma}{Lemma}[section]
\newtheorem{corollary}{Corollary}[section]

\theoremstyle{definition}
\newtheorem{definition}{Definition}[section]
\newtheorem{example}{Example}[section]

\theoremstyle{remark}
\newtheorem{remark}{Remark}[section]

\crefname{theorem}{theorem}{theorems}
\crefname{proposition}{proposition}{propositions}
\crefname{lemma}{lemma}{lemmas}
\crefname{corollary}{corollary}{corollaries}
\crefname{definition}{definition}{definitions}
\crefname{remark}{remark}{remarks}
\crefname{example}{example}{examples}


\title{Named Grid Covers under Row--Column Margins:\\
Three Exactness Levels and Largest-Piece Localization}
\author{Oleksiy Babanskyy}
\date{}

\begin{document}

\maketitle

\begin{abstract}
We study finite named grid-piece placements whose aggregate row and column
margins agree with a target, while overlaps and uncovered cells are allowed
inside the relaxed fiber.  Exact covers form a distinguished sublocus.  This
gives three different questions: whether a margin fiber supports an exact
cover, whether every margin realization is exact, and whether every
realization can reach an exact cover by replacing exactly two named pieces
at a time.  After fixing a pose of a largest selected piece, complement area
\(q\) confines every participating residual placement to at most \(q\)
active rows, \(q\) active columns, and \(q^2\) cells.  We derive fixed-
parameter algorithms for all three questions in \((q,h)\), where \(h\) bounds
the orientation group, and the explicit fiber bound
\[
  h(q+1)(h q^2)^q=2^{O(q\log q)}
\]
for fixed \(h\).  Source-locked finite libraries show that the three
predicates are genuinely distinct.  A bounded 72-variable example further
separates connectivity of one nonnegative fiber, generation of the integer
relation lattice, and Markov connectivity over all right-hand sides.
\end{abstract}

\begin{center}
\small\emph{Author manuscript. See the repository claim boundary and reproducibility package.}
\end{center}

\tableofcontents

\section{Introduction}\label{sec:introduction}

Tomographic reconstruction asks what can be recovered from projections rather
than from direct cell data.  In grid tilings the feasible objects are usually
required to be disjoint exact tilings, possibly with typed projections or a
fixed tile shape.  This literature already contains reconstruction
algorithms, hardness results, exact domino models, and heterogeneous or
polyatomic formulations
\cite{HermanKuba1999,ChrobakEtAl2003,ChrobakDurr2001,DurrEtAl2000,
FrosiniSimi2004,FrosiniSimi2005,GritzmannLangfeld2020,
ChrobakEtAl2012}.  We ask a different, deliberately narrow question.
Suppose that one named copy of each selected piece is placed inside a target
and only the \emph{aggregate} row and column margins are required to match.
Overlaps and holes are allowed at this relaxed stage.  What does that margin
condition say about exact covers, and how hard is it to recognize the
possible failure modes?

The distinction matters because a correct aggregate margin vector need not
certify exact coverage cell by cell.  We separate three predicates:

\begin{enumerate}
  \item \emph{support}: a nonempty margin fiber contains an exact cover;
  \item \emph{purity}: every point of the margin fiber is an exact cover;
  \item \emph{repair}: every fiber component, under replacement of exactly
        two named placements, contains an exact cover.
\end{enumerate}

Purity implies repair, and repair implies support, but neither implication
reverses.  The strictness is witnessed by complete finite certificates for
three source-locked named libraries.  These witnesses organize the problem;
the logical implications themselves are elementary, and general
reconfiguration from infeasible states is prior art
\cite{FellowsEtAl2023}.

Our main uniform result uses the \emph{largest-piece residual area}
\[
  q=|T|-|L|,
\]
where \(L\) is a largest selected piece.  Once a legal pose of \(L\) is
fixed, the residual row and column margins have total mass \(q\).
Consequently every residual placement that can occur in the margin fiber is
confined to at most \(q\) active rows, at most \(q\) active columns, and
hence at most \(q^2\) active cells.  A translation bound then gives at most
\(h(q+1)\) poses for the largest piece and yields a
\(2^{O(q\log q)}\) bound on the complete margin fiber for fixed orientation
bound \(h\).  Enumerating that bounded fiber decides support, purity, and
exact-two repair.

\subsection{Contributions and scope}

The paper makes the following precise statements.

\begin{itemize}
  \item We define the named one-copy aggregate-margin fiber and its exact-cover
        sublocus, together with support, purity, and exact-two repair.
  \item We prove the active-carrier lemma and fixed-parameter tractability of
        all three predicates in \((q,h)\) for finite cell-list pieces.
  \item We give an exact participation-type compression for existential
        decision and compressed classification in a finite named library,
        under an explicit residual-area threshold.  Literal listing remains
        output-sensitive.
  \item We record a simple overlap energy whose strict descent is sufficient
        for repair, without claiming necessity.
  \item We use certified C32, C46, and P49 libraries to separate the three
        exactness levels.
  \item In one frozen P21 configuration we distinguish connectivity of a
        selected nonnegative fiber, integral relation-lattice generation, and
        all-right-hand-side Markov connectivity.
\end{itemize}

The only uniform mathematical headline is the active-carrier/FPT theorem.
We do not claim to introduce heterogeneous tiling tomography, pair
replacement, infeasible-state repair, or algebraic tiling connectivity.
Locked pair phenomena and binomial move connectivity for exact tilings have
already been studied \cite{TuckerFoltz2023,GrossYamzon2021}.  The algebraic
language in the P21 section is classical
\cite{DiaconisSturmfels1998,AokiTakemuraYoshida2008,
CharalambousThomaVladoiu2016}.  No priority or ``first FPT'' claim is made.

\subsection{Organization}

\Cref{sec:model} fixes the finite input contract and the exact-two graph.
\Cref{sec:exactness,sec:libraries} establish the hierarchy and its finite
witnesses.  \Cref{sec:localization,sec:fpt} prove localization and the FPT
theorem.  \Cref{sec:types-energy} treats implicit libraries and the overlap
potential.  \Cref{sec:p21} gives the bounded algebraic case study.
\Cref{sec:limits} states the stop boundary.  The appendices record the
evidence and reproduction contracts.

\section{Named exact covers and tomographic fibers}\label{sec:model}

\subsection{Finite placement model}

Let \(T\subset\ZZ^2\) be a finite target cell set.  A library consists of
nonempty, named, finite cell-list pieces.  Legal poses are obtained from a
supplied finite orientation group \(H\), with \(|H|\leq h\), followed by
integer translations; every legal pose must be contained in \(T\).
Explicitly, each \(g\in H\) acts by a lattice automorphism
\(g:\ZZ^2\to\ZZ^2\), the input supplies the action (or names one fixed
action), and the pose catalogue consists of the distinct cell masks
\(gP+t\subseteq T\).  Repeated masks produced by symmetries are identified.
Names remain distinct even if two pieces have the same shape or area.

Fix a selected named set
\[
  E=\{P_0,\ldots,P_r\},\qquad \sum_{i=0}^r |P_i|=|T|.
\]
The area equality is part of the decision domain.  Area-incompatible
selections are rejected during preprocessing.  Write \(\PP_i\) for the
finite catalogue of legal poses of \(P_i\).

For a cell set \(A\subseteq T\), let \(X(A)\) be its joint row-and-column
margin vector.  A one-copy placement tuple is
\[
  x=(A_0,\ldots,A_r)\in\prod_{i=0}^r\PP_i.
\]
Its cell coverage is
\[
  c_x(u)=|\{i:u\in A_i\}|,\qquad u\in T.
\]

\begin{definition}[Margin fiber and exact-cover sublocus]\label{def:fiber}
For the target margin \(X(T)\), define
\[
  \FFib_E=\left\{x:\sum_i X(A_i)=X(T)\right\}
\]
and
\[
  \TT_E=\{x\in\FFib_E:c_x=\one_T\}.
\]
Thus \(\TT_E\subseteq\FFib_E\).  Points of \(\FFib_E\setminus\TT_E\)
may have overlaps and uncovered cells, but every chosen pose still lies in
the target.
\end{definition}

The row and column margins are aggregate: the contribution of a piece name
is not separately observed.  This differs from typed projections that record
one margin vector per tile type.  It also differs from exact tiling models in
which disjointness is imposed before the projection constraints.

\subsection{The replacement graph}

\begin{definition}[Exact-two replacement graph]\label{def:graph}
The graph \(\GG_E^{(2)}\) has vertex set \(\FFib_E\).  Two distinct tuples
are adjacent precisely when they differ in exactly two named coordinates.
\end{definition}

An edge is one graph step, although its endpoints have Hamming distance two
as named placement tuples.  We also use the explicitly labelled variant
\(\GG_E^{(\leq2)}\), which adds edges between tuples differing in one named
coordinate.  These graphs need not agree.  Every unqualified repair
statement in this paper refers to \(\GG_E^{(2)}\).

\begin{definition}[Three exactness predicates]\label{def:predicates}
For every area-compatible selected set \(E\), define the following three
Boolean predicates:
\begin{itemize}
  \item \emph{support-exact} if
        \(\FFib_E\neq\varnothing\) exactly when
        \(\TT_E\neq\varnothing\);
  \item \emph{fiber-pure} if \(\FFib_E=\TT_E\);
  \item \emph{repair-exact} if every connected component of
        \(\GG_E^{(2)}\) meets \(\TT_E\).
\end{itemize}
If repair-exactness holds, the repair radius is
\[
  \rho(E)=\max_{x\in\FFib_E}\dist_{\GG_E^{(2)}}(x,\TT_E).
\]
\end{definition}

Equivalently, a finite library has any one of these properties on a declared
selection domain \(\mathcal D\) when the corresponding predicate holds for
every \(E\in\mathcal D\).  Finite-library tables in this paper always state
their domain explicitly.  Empty fibers satisfy all three predicates
vacuously and have radius zero.
If a component is disjoint from \(\TT_E\), its vertices are called
\emph{trapped}; they are not assigned a finite repair distance.

\subsection{Input-size contract}

The target and piece masks are explicit finite cell lists.  The orientation
group is supplied explicitly by its lattice action, or by a fixed named
lattice action of size at most \(h\).
Translations and margin vectors use binary-encoded integer coordinates.
The input length \(|I|\) includes these lists and coordinates.  The theorem
does not address arbitrary continuous bodies, oracle-described shapes,
unlimited tile multiplicities, or anonymous tile types.

\section{Three levels of tomographic exactness}\label{sec:exactness}

\begin{proposition}[Exactness hierarchy]\label{prop:hierarchy}
For every area-compatible selected set,
\[
  \text{fiber purity}\ \Longrightarrow\
  \text{repair-exactness}\ \Longrightarrow\
  \text{support-exactness}.
\]
\end{proposition}

\begin{proof}
If \(\FFib_E=\TT_E\), every graph vertex is already an exact cover, so every
component meets \(\TT_E\).  If repair-exactness holds and the fiber is
nonempty, choose any component; it contains a point of \(\TT_E\).  The reverse
implication \(\TT_E\neq\varnothing\Rightarrow\FFib_E\neq\varnothing\)
always holds because \(\TT_E\subseteq\FFib_E\).  The empty case follows from
the conventions in \cref{def:predicates}.
\end{proof}

The value of \cref{prop:hierarchy} is organizational rather than logical
difficulty.  Generic discovery of a feasible solution from an infeasible
state under restricted modifications belongs to the wider reconfiguration
literature \cite{FellowsEtAl2023}.  Here the predicates are tied to one
specific aggregate tomographic relaxation and one specific named
replacement graph.

\begin{proposition}[Strictness in locked finite libraries]
\label{prop:strictness}
Neither implication in \cref{prop:hierarchy} reverses within the certified
finite named-library class:
\begin{enumerate}
  \item C46 is repair-exact with radius one but is not fiber-pure.
  \item The P49 row of zero-based indices \((4,9,10,11)\) is support-exact,
        but eight non-tiling nodes are trapped singleton components.
  \item In a separate failure, the P49 row \((0,1,3,14)\) has a nonempty
        margin fiber and no exact cover.
\end{enumerate}
C32 is a positive pure anchor.
\end{proposition}

\begin{proof}
All assertions follow from the exhaustive finite proposition in
\cref{prop:finite-panel}; the certificate and independent replay contract is
recorded in \cref{app:evidence}.
\end{proof}

These examples prove non-collapse of the predicates in the locked model.
They do not classify shapes, characterize repairability, or imply that
impure fibers normally repair.

\section{Certified foundational libraries}\label{sec:libraries}

\subsection{Evidence layers and label conventions}

Each finite statement below is separated into four layers.

\begin{enumerate}
  \item A source lock records the asset descriptor, signature, direct locator,
        access date, and SHA-256 digest.  The rasters are not redistributed.
  \item A discrete normalization records the target cells, named piece masks,
        coordinate convention, and full \(D_4\) pose convention.
  \item An exact-cover census gives all tilings for every positive
        area-compatible row and exhaustive zero certificates for the rest.
  \item An independent tomographic replay reconstructs every margin fiber,
        exact-cover subset, and exact-two graph downstream of the frozen
        normalized transcription.
\end{enumerate}

The finite libraries C32, C46, and P49 are normalized from \(4\times4\)
geomagic specimens created and published by Lee Sallows
\cite{Sallows2011,SallowsGallery}.  The piece geometry, targets, and displayed
arrangements are Sallows' source material; Peter Cameron's finite
group-action formulation is prior conceptual context \cite{Cameron2011}.
The present project's contribution begins with the source-locked discrete
normalization and consists of exhaustive pose and exact-cover censuses,
aggregate-margin fibers, replacement graphs, and the uniform
largest-piece-localization theorem.  We do not claim authorship of the
underlying geomagic specimens.  The 86-row additive carrier shared by C32
and C46 is imported and is not presented as a new additive theorem.

C32 and C46 have distinct piece areas \(1,\ldots,16\), so a selected row lists
both labels and areas.  P49 has repeated areas; its row tuples always list
zero-based locked piece indices.  In particular,
\((4,9,10,11)\) has areas \(\{4,13,1,14\}\), whereas
\((0,1,3,14)\) has areas \(\{15,2,3,12\}\).

\subsection{Source locks}

C32 is the Sallows gallery asset \emph{4x4 Durer5}, signed LS 11-06, with locked SHA-256
\[
\mathtt{c6093ccd58e725d6efde1750cccf47124a56d6a335fbe3c9b70c8ef66a9bf2cb}.
\]
Its target is a \(6\times6\) square with zero-based holes \((1,1)\) and
\((4,4)\).  Of twenty displayed panels, twelve already use these holes and
eight use the horizontal reflection.  Reflecting the latter group into the
canonical target gives sixteen distinct labelled tilings.

C46 is the Sallows gallery asset \emph{4x4 most perfect (2b)}, signed LS
4-09, with locked SHA-256
\[
\mathtt{86430abfcfe5793fecdedf42a84a0e4f08e16636e5691ab25cf4509e9c69ef00}.
\]
It uses the same target and area-labelled piece convention.  Its 36 displayed
panels agree on the target mask and on the recovered connected color
components.

P49 is the source-locked Sallows specimen \emph{Art of Fugue}, signed LS
7-09, with locked SHA-256
\[
\mathtt{77efa6d693726aa026fce49e14480b933606ce552ff51a24eca453626d175c1e}.
\]
Its target is a \(6\times6\) square with the four corners removed.  Its
zero-based area sequence is
\[
 (15,2,12,3,4,11,7,10,4,13,1,14,9,6,12,5).
\]

\begin{proposition}[Complete finite panel]\label{prop:finite-panel}
Under the locked models and the semantics of \cref{sec:model}, the complete
area-compatible carriers have the following exact behavior.

\begin{center}
\small
\begin{tabular}{@{}lrrrrlll@{}}
\toprule
Library & Rows & Tile/zero & Margin & Fiber \(=\) tile+extra
& Support & Pure & Repair\\
\midrule
C32 & 86 & \(20/66\) & 20 & \(136=136+0\)
& yes & yes & yes, \(\rho=0\)\\
C46 & 86 & \(44/42\) & 44 & \(352=344+8\)
& yes & no & yes, \(\rho=1\)\\
P49 & 94 & \(60/34\) & 61 & \(1344=1208+136\)
& no & no & no\\
\bottomrule
\end{tabular}
\end{center}
\end{proposition}

Here \emph{Margin} is the number of nonempty margin fibers, while
\emph{Tile/zero} and the displayed fiber decomposition report the corresponding
tiling counts and exhaustive zero certificates.  Thus no positive exactness
label is stated without accompanying existence data.

\begin{proof}
For each locked catalogue, the primary enumerator exhausts every
area-compatible row and stores complete labelled exact covers or a zero
certificate.  A separately implemented replay regenerates \(D_4\) poses,
constructs the fibers by direct \(2+2\) meet-in-the-middle matching, and
agrees on every row count, edge count, component, radius, and verdict.
Positive and negative controls pass, while the registered adversarial
mutations are rejected.  Exact commands and artifact paths appear in
\cref{app:evidence,app:reproduction}.
\end{proof}

The replay is an independent downstream enumeration from these frozen
normalized transcriptions.  Because the original source rasters are not
redistributed, the reviewer package does not independently reproduce image
acquisition, segmentation, or transcription from Sallows' raster bytes.  It
verifies the normalized model and every mathematical deduction downstream of
that source lock.

\subsection{Witness rows}

C32 has 136 margin nodes and every one is an exact cover.  It is therefore
pure.  Seven of its positive rows are non-affine; this is useful as a warning
against an affine-selection characterization but is not needed below.

C46 has 344 tilings and eight non-tiling margin matches.  All eight extras
occur in two rows:

\begin{center}
\small
\begin{tabular}{@{}lrrrrrr@{}}
\toprule
Row & Fiber & Tilings & Extras & Components & Edges & Radius\\
\midrule
\(\{1,3,14,16\}\) & 20 & 16 & 4 & 2 & 22 & 1\\
\(\{2,3,13,16\}\) & 8 & 4 & 4 & 4 & 4 & 1\\
\bottomrule
\end{tabular}
\end{center}

Every extra is adjacent to a tiling.  Hence C46 is repair-exact but not pure.
The eight positive rows not displayed in the source are described only as
\emph{not source-displayed}: the source never claimed that its panel list was
exhaustive.

P49 has two independent failure modes:

\begin{center}
\small
\begin{tabular}{@{}lrrrrrr@{}}
\toprule
Indices & Fiber & Tilings & Extras & Components & Edges
& Trapped comps./nodes\\
\midrule
\((4,9,10,11)\) & 24 & 8 & 16 & 16 & 8 & \(8/8\)\\
\((0,1,3,14)\) & 8 & 0 & 8 & 8 & 0 & \(8/8\)\\
\bottomrule
\end{tabular}
\end{center}

In the first row, eight extras pair with tilings and eight extras are
\emph{eight separate singleton components}.  In the second row, all eight
nodes are trapped singletons and no tiling exists.  Across P49 there are
sixteen trapped components.  This corrects the inaccurate shorthand of one
eight-node trapped component.

\section{Largest-piece residual localization}\label{sec:localization}

Choose a largest selected name \(L\), breaking area ties by a fixed name
order, and put
\[
  q=|T|-|L|.
\]
The residual pieces are nonempty and their total area is \(q\), so there are
at most \(q\) of them.

\subsection{A translation cap}

For finite nonempty \(A,B\subset\ZZ^2\), define
\[
  D(A,B)=\{t\in\ZZ^2:A+t\subseteq B\}.
\]

\begin{lemma}[Translation cap]\label{lem:translation}
For finite nonempty \(A,B\subset\ZZ^2\),
\[
 |D(A,B)|\leq \max(0,|B|-|A|+1).
\]
Equivalently, if \(D(A,B)\) is nonempty, then
\[
 |D(A,B)|\leq |B|-|A|+1.
\]
\end{lemma}

\begin{proof}
The empty case is immediate.  Suppose \(D=D(A,B)\) is nonempty.  Then
\(A+D\subseteq B\).  Choose an integer linear functional \(\ell\) injective
on the finite set \(A+D\); avoiding the finitely many forbidden integer
slopes suffices.  It is then injective on each factor as well.  If
\(\ell(A)=\{a_1<\cdots<a_m\}\) and
\(\ell(D)=\{d_1<\cdots<d_n\}\), the chain
\[
 a_1+d_1<\cdots<a_m+d_1<a_m+d_2<\cdots<a_m+d_n
\]
contains \(m+n-1\) distinct sums.  Therefore
\[
 |B|\geq|A+D|=|\ell(A)+\ell(D)|
 \geq |A|+|D|-1,
\]
and rearrangement proves the claim.
\end{proof}

The sumset mechanism is classical; the lemma is included to make the
algorithmic constants explicit.  The maximum in the unconditional statement
is essential when no translation exists.

\subsection{The active carrier}

Freeze one legal pose \(A\) of \(L\) and let
\[
  b=X(T)-X(A).
\]
Since \(A\subseteq T\), the vector \(b\) is nonnegative.  Its row coordinates
sum to \(q\), and its column coordinates also sum to \(q\).

\begin{lemma}[Active carrier]\label{lem:active}
After a pose \(A\) of \(L\) is fixed, every residual placement participating
in a target-margin solution is confined to at most \(q\) active rows, at
most \(q\) active columns, and therefore at most \(q^2\) target cells.
\end{lemma}

\begin{proof}
At most \(q\) row coordinates and at most \(q\) column coordinates of \(b\)
are positive.  If a residual pose \(Q\) occurs in a tuple whose residual
margins sum to \(b\), nonnegativity gives \(X(Q)\leq b\) coordinatewise.
Thus \(Q\) uses only rows and columns where \(b\) is positive.  It is
contained in
\[
  G_A=T\cap(\text{active rows}\times\text{active columns}),
  \qquad |G_A|\leq q^2.
\]
\end{proof}

The proof uses both containment of every pose in \(T\) and the declared
row/column finite-grid margins.  It is not a statement about continuous
tomography.

\begin{corollary}[Pose bounds]\label{cor:pose-bounds}
The largest piece has at most \(h(q+1)\) legal poses.  After one is fixed,
each residual piece has at most \(hq^2\) margin-compatible candidate poses.
\end{corollary}

\begin{proof}
Apply \cref{lem:translation} to each oriented copy of \(L\) in \(T\):
every nonempty translation set has size at most
\(|T|-|L|+1=q+1\).  Sum over at most \(h\) orientations.  Any residual pose
compatible with the residual margin vector after \(A\) is fixed is a nonempty
subset of the carrier \(G_A\).  The same lemma gives at most \(q^2\)
translations per orientation, and hence at most \(hq^2\) candidate poses.
This catalogue may overcount poses that do not extend to a full fiber tuple.
\end{proof}

\section{Fixed-parameter decision and radius}\label{sec:fpt}

\subsection{Bounded dynamic programs}

Fix a pose \(A\) of the canonical largest piece \(L\), and retain the
notation \(b=X(T)-X(A)\).  For margin feasibility, a partial-placement
dynamic program needs only residual margin states bounded coordinatewise by
\(b\).  Its exact state box has size
\[
  B(b)=\prod_j(b_j+1).
\]
For nonnegative integers \(n\), \(n+1\leq2^n\).  The row coordinates of
\(b\) sum to \(q\), as do the column coordinates, so
\[
  B(b)\leq 2^q2^q=4^q.
\]
For exact coverage, the residual target \(T\setminus A\) has \(q\) cells.
A subset-mask dynamic program therefore has at most \(2^q\) states.
Disjoining over the poses of \(L\) decides whether the margin fiber and the
exact-cover sublocus are empty.

To decide purity and repair, we reconstruct all accepting margin paths.  This
is safe because the entire fiber is itself parameter-bounded.

\begin{theorem}[Largest-piece residual-area FPT theorem]\label{thm:fpt}
For an area-compatible fixed selected set of named finite grid pieces and a
finite orientation group of size at most \(h\), support-exactness, fiber
purity, and repair-exactness in the exact-two graph are fixed-parameter
tractable in \((q,h)\).  The full margin fiber satisfies
\[
 |\FFib_E|\leq h(q+1)(h q^2)^q
              =2^{O(q\log(qh))}.
\]
For fixed \(h\), this simplifies to \(2^{O(q\log q)}\).  The replacement
graph can be constructed in
\[
 2^{O(q\log(qh))}\operatorname{poly}(|I|)
\]
time, and hence in \(2^{O(q\log q)}\operatorname{poly}(|I|)\) time when
\(h\) is fixed.
\end{theorem}

\begin{proof}
By \cref{cor:pose-bounds}, \(L\) has at most \(h(q+1)\) legal poses.  For a
fixed largest pose, there are at most \(q\) residual pieces and each has at
most \(hq^2\) margin-compatible candidate poses.  Hence
\[
 |\FFib_E|
 \leq h(q+1)(hq^2)^q.
\]
This includes nonaccepting combinations in its upper estimate and is
therefore valid for the fiber.  The case \(q=0\) is handled directly.

Enumerating the parameter-bounded fiber and testing the cell coverage of
each tuple decides purity.  Support is decided by comparing nonemptiness of
\(\FFib_E\) and \(\TT_E\).

Hash all fiber tuples.  For each tuple, choose two named coordinates,
enumerate alternatives from their candidate-pose catalogues, and retain
the resulting tuple precisely when it occurs in the hash table.  There are
at most polynomially many choices outside the parameter-bounded pose
factors, so this constructs \(\GG_E^{(2)}\) in
\(2^{O(q\log(qh))}\operatorname{poly}(|I|)\) time.  Mark its connected
components and test whether each intersects \(\TT_E\).  This decides
repair-exactness.
\end{proof}

\begin{remark}
Adding the one-name cases constructs \(\GG_E^{(\leq2)}\) in the same FPT
class.  This does not identify the two graphs or transfer finite component
counts between them.
\end{remark}

\begin{corollary}[Existence bound for repair radius]\label{cor:radius}
Every repair-exact fiber satisfies
\[
 \rho(E)\leq|\FFib_E|-1
 \leq h(q+1)(h q^2)^q-1.
\]
\end{corollary}

\begin{proof}
In a component meeting \(\TT_E\), choose a shortest path from a vertex to a
tiling.  A shortest path is simple and hence has at most the component size
minus one edges.
\end{proof}

The bound in \cref{cor:radius} is an existence bound.  It is not asserted to
be polynomial, linear, sharp, or representative of observed runtime.

\section{Implicit libraries and a sufficient repair potential}
\label{sec:types-energy}

\subsection{Exact participation types}

\Cref{thm:fpt} concerns a fixed selected named set.  To avoid conflating that
problem with a search through a large library, fix the following separate
parameterized problem.  The input is a finite target \(T\), a finite named
cell-list library \(\mathcal Q\), the explicit orientation action, an integer
\(Q\), and one of the three predicates in \cref{def:predicates}.  The question
is whether there is an area-compatible selected set \(E\subseteq\mathcal Q\)
whose canonical largest name \(L(E)\) satisfies
\[
 |T|-|L(E)|\leq Q
\]
and whose margin fiber has the requested predicate.  The classification
version returns one compressed record per behavior class; it does not list
all named selected sets.

Fix a candidate largest name \(L\), order its legal poses as
\(\mathcal A_L=(A_1,\ldots,A_s)\), and put \(q_L=|T|-|L|\).  A residual name
\(P\) is \emph{eligible for \(L\)} when \(|P|<|L|\), or when
\(|P|=|L|\) and the fixed tie order chooses \(L\) before \(P\).  For
\(A_j\in\mathcal A_L\), define
\[
 b_j=X(T)-X(A_j),\qquad
 \mathcal M_j(P)=\{B\in\PP_P:X(B)\leq b_j\text{ coordinatewise}\}.
\]
Every mask in \(\mathcal M_j(P)\) lies in the active carrier \(G_{A_j}\).
Use the common labelled carrier
\[
 S_L=\bigcup_{j=1}^{s}G_{A_j},
 \qquad |S_L|\leq h(Q+1)Q^2.
\]

\begin{definition}[Exact participation type]\label{def:participation-type}
The exact participation type of a residual name \(P\), relative to \(L\), is
\[
 \tau_L(P)=
 \bigl(|P|,\text{eligibility},
       (\mathcal M_1(P),\ldots,\mathcal M_s(P))\bigr).
\]
Each placement is stored as its actual labelled cell mask in \(S_L\), not as
an index local to one scenario.  Thus the record retains the identity of a
placement mask that occurs for more than one largest-piece pose.  Let
\(m_L(\tau)\) be the number of available residual names of type \(\tau\).
\end{definition}

\begin{lemma}[Type-isomorphism lemma]\label{lem:type-isomorphism}
Fix \(L\).  Suppose two residual named multisets
\((P_1,\ldots,P_r)\) and \((P'_1,\ldots,P'_r)\) admit a bijection \(\sigma\)
with \(\tau_L(P_i)=\tau_L(P'_{\sigma(i)})\) for every \(i\).  Replacing each
name by its matched name and retaining the same labelled placement mask gives
an isomorphism of margin fibers.  This isomorphism maps exact covers to exact
covers and is a graph isomorphism for both \(\GG^{(2)}\) and
\(\GG^{(\leq2)}\).
\end{lemma}

\begin{proof}
For every ordered pose \(A_j\) of \(L\), equality of types gives exactly the
same allowed cell masks in every matched residual coordinate.  Hence the
coordinatewise renaming is a bijection of placement tuples and preserves the
sum of margin vectors.  Because the actual masks are retained, it also
preserves the coverage function and therefore the exact-cover sublocus.
Two tuples differ in a named coordinate before the renaming exactly when
their images differ in the matched coordinate.  This remains true when the
largest pose changes: cross-scenario mask identity records which unchanged
residual coordinates are literally the same placement.  Hamming distance
one or two, and therefore both replacement graphs, are preserved.
\end{proof}

\begin{proposition}[Compressed finite-library decision]\label{prop:types}
The existential problem above and its compressed behavior classification are
FPT in \((Q,h)\).  For a candidate \(L\), a compressed record contains \(L\),
an exact participation-type count vector \(k=(k_\tau)_\tau\), the number
\[
 \prod_\tau \binom{m_L(\tau)}{k_\tau}
\]
of named residual selections represented by that vector, and the resulting
fiber predicates and counts.
\end{proposition}

\begin{proof}
For each candidate \(L\) with \(0\leq q_L\leq Q\), the ordered list of
largest poses and the carrier \(S_L\) have size bounded by a function of
\((Q,h)\).  Area, eligibility, and the exact mask family for every ordered
largest pose therefore yield only \(f(Q,h)\) possible types.  The selected
residual names are nonempty and have total area \(q_L\), so at most \(Q\) of
them are chosen.  We may truncate the available multiplicity of every type
at \(Q\), enumerate the type-count vectors satisfying
\[
 \sum_\tau k_\tau\,\operatorname{area}(\tau)=q_L
\]
and the eligibility rule, and run \cref{thm:fpt} on one representative of
each vector.  The type-isomorphism lemma proves that this representative has
exactly the same fiber, exact-cover sublocus, and move graphs as every named
selection counted by the displayed product.  Constructing the signatures
for all candidate largest names is polynomial in the explicit input size;
all superpolynomial enumeration is confined to a function of \((Q,h)\).
\end{proof}

\begin{example}[Why literal listing is output-sensitive]
Take a \(1\times2Q\) target, one named \(Q\)-bar, and \(N\) named monominoes.
There are \(\binom NQ\) distinct valid named selections.  Consequently the
literal list of all selections cannot, in general, be produced in total FPT
time independent of output size.  No Delay-FPT enumerator is claimed here;
the distinction follows the standard parameterized-enumeration vocabulary
\cite{CreignouEtAl2017}.
\end{example}
\subsection{Overlap energy}

For \(x\in\FFib_E\), define
\[
 \Phi(x)=\sum_{u\in T}\max(c_x(u)-1,0).
\]

\begin{proposition}[Energy identity and strict descent]\label{prop:energy}
For every margin node \(x\),
\[
 \Phi(x)
 =|T|-\left|\bigcup_i A_i\right|
 =\frac12\|c_x-\one_T\|_1,
 \qquad 0\leq\Phi(x)\leq q.
\]
Moreover, \(\Phi(x)=0\) exactly when \(x\in\TT_E\).  If every non-tiling node
has an exact-two neighbor of strictly smaller energy, then the fiber is
repair-exact and \(\rho(E)\leq q\).
\end{proposition}

\begin{proof}
The chosen pose areas sum to \(|T|\), so
\[
 \sum_u c_x(u)=|T|.
\]
The overlap mass \(\sum_u\max(c_x(u)-1,0)\) is therefore the total area minus
the size of the covered union.  Also
\(\sum_u(c_x(u)-1)=0\): positive deviation equals uncovered mass, proving the
\(\ell_1\) identity.  The largest pose covers \(|T|-q\) distinct cells, so at
most \(q\) cells can be missing from the union, whence \(0\leq\Phi\leq q\).
Energy zero is equivalent to unit coverage everywhere.  Finally, repeated
strict descent in the nonnegative integer \(\Phi\) reaches zero in at most
\(\Phi(x)\leq q\) steps.
\end{proof}

Strict descent is sufficient, not necessary.  Certified P21 and P26 rows
contain repairable nonzero local minima that require a nondecreasing or
uphill step.  Thus local minima do not characterize trapped components.

\section{P21: three non-equivalent connectivity questions}\label{sec:p21}

The uniform theorem concerns one selected one-copy margin fiber.  The P21
example is included to show what that fiber cannot certify.  It distinguishes:
\begin{enumerate}
  \item connectivity of one fixed nonnegative fiber;
  \item generation of the complete integer relation lattice;
  \item connectivity of every nonnegative fiber, equivalently the Markov
        basis requirement.
\end{enumerate}
The first fixes one right-hand side.  The second allows signed intermediate
combinations.  The third varies the right-hand side and must preserve
nonnegativity throughout \cite{DiaconisSturmfels1998}.  Algebraic
formulations of grid-tiling flip ideals also precede this example
\cite{Chin2019,GrossYamzon2021}; the claims below are only for the frozen
72-variable configuration.

\subsection{Frozen configuration}

The source-locked specimen is \emph{4x4 square target}, signed LS 11-02, with
SHA-256
\[
\mathtt{2248f6e9d0d4cb793ef2ff8069a37ae62ed67629cee6263a0f7071ae4a274120}.
\]
Its target is the \(5\times5\) square.  We freeze the zero-based piece-index
row
\[
 E=(4,5,12,13),
\]
whose areas are \((7,6,6,6)\).  Only placements participating in the selected
margin fiber are retained.  Their block counts are \(8,24,16,24\), for 72
variables.

For each participating placement \(P\) in named block \(i\), the
configuration matrix \(A\) has column
\[
 e_i\mid X_{\rm row}(P)\mid X_{\rm col}(P).
\]
Thus \(A\) is a \(14\times72\) integer matrix of rank 12 and
\[
 K=\ker_{\ZZ}(A)
\]
has rank 60.

For \(s=1,2,3,4\), let \(M_{\leq s}\) be the integer span of all squarefree
one-copy relations whose two placement tuples differ in at most \(s\) named
piece positions.  This \emph{piece support} is not toric degree, Graver
support, or a claim of Graver primitiveness.

\subsection{The fixed 32-node fiber}

The chosen fiber contains 32 nodes: 8 exact tilings and 24 non-tiling margin
matches.  Its complete support filtration is:

\begin{center}
\small
\begin{tabular}{@{}lrrrrl@{}}
\toprule
Family & New rels. & Lattice rank & \(\rank K/M\) & Components & Verdict\\
\midrule
\(M_{\leq2}\) & 144 & 34 & 26 & 16 & 16 trapped nodes\\
\(M_{\leq3}\) & 932 & 51 & 9 & 4 & all reach tilings\\
\(M_{\leq3}+\langle O_{714}\rangle\) & one orbit & 54 & 6 & 1 & connected\\
\(M_{\leq4}\) & 7,444 & 60 & 0 & 1 & complete graph\\
\bottomrule
\end{tabular}
\end{center}

The relation counts are new at the indicated exact support.  The complete
\(M_{\leq3}\) and \(M_{\leq4}\) catalogues contain 1,076 and 8,520 relations,
respectively.  At support at most two there are sixteen two-node components:
eight pair a tiling with an extra, while eight contain the sixteen trapped
nodes.  At support three there are four eight-node components, but every node
reaches a tiling.  At support four the graph is complete.

\begin{proposition}[One connected fiber does not generate the lattice]
\label{prop:p21-one}
The family \(M_{\leq3}+\langle O_{714}\rangle\) connects the fixed 32-node
fiber, but has lattice rank \(54<60=\rank K\).
\end{proposition}

\begin{proof}
The order-eight target action partitions the 7,444 support-four relations
into 1,119 complete orbits.  Exactly 36 single orbit modules connect the four
components of the support-three graph.  Direct Hermite/Smith reduction gives
rank 54, with torsion-free quotient rank six, for every one of these choices;
\(O_{714}\) is the canonical stored witness.
\end{proof}

\subsection{The missing integral channels}

The quotient \(Q=K/M_{\leq3}\) is torsion-free of rank nine.  Among the 1,119
support-four orbits, 87 cyclic orbit modules have quotient rank three.  They
fall into three exact integral families of sizes
\[
 6,\quad3,\quad78.
\]
Exactly three complete target-action orbits are necessary and sufficient to
generate \(Q\) integrally.  Every minimum triple chooses one orbit from each
family, giving \(6\cdot3\cdot78=1,404\) minimum triples.  The stored canonical
triple is \((283,603,714)\), with determinant \(-1\) in the quotient basis.

The fixed 32-node fiber activates \(0,0,36\) orbits in these three families.
It is blind to the first two necessary integral channels.  These numbers are
configuration-specific finite facts, not a general equivariant-basis
theorem; compare the general lifting context in
\cite{RauhSullivant2019}.

\subsection{Lattice generation is not Markov connectivity}

All 60 Smith factors for \(M_{\leq4}\subseteq K\) are one, hence
\[
 M_{\leq4}=K
\]
integrally.  This permits signed decompositions of every relation.  It does
not imply that a decomposition can be ordered while staying nonnegative.

\begin{proposition}[All-right-hand-side no-go]\label{prop:p21-markov}
Although \(M_{\leq4}=K\), the squarefree one-copy support-at-most-four family
is not a Markov basis for the frozen 72-variable matrix.
\end{proposition}

\begin{proof}
The complete quadratic census has 218 indispensable binomials: 144
cross-block quadrics occur in the one-copy family, while 74 within-block
quadrics do not.  A canonical missing quadratic lies in a two-monomial fiber
with block count \((2,0,0,0)\).  No one-copy cross-block move applies in that
fiber, so it is disconnected.  Indispensability is certified using the
standard binomial-fiber criterion
\cite{AokiTakemuraYoshida2008,CharalambousThomaVladoiu2016}.
\end{proof}

A separate two-monomial fiber with block count \((0,0,0,7)\) forces an
indispensable binomial of toric degree seven.  Therefore every Markov basis
has degree \emph{at least} seven.  The bounded minimum-generator census is:

\begin{center}
\begin{tabular}{@{}rrr@{}}
\toprule
Toric degree & Betti fibers & Minimum generators forced\\
\midrule
2 & 218 & 218\\
3 & 996 & 996\\
4 & 5,705 & 5,727\\
\bottomrule
\end{tabular}
\end{center}

At least 6,941 minimal generators are forced through degree four.  No upper
bound of seven, exact Markov degree, complete basis, Graver basis, universal
Markov basis, or blockwise lifting theorem follows.

\section{Limits and open directions}\label{sec:limits}

The theorem closes the decision problem under a precise finite contract; it
does not close the wider geometry or algebra.  The following directions are
open and logically independent of the present manuscript.

\paragraph{Sharper radius.}
\Cref{cor:radius} bounds the radius by the number of fiber states minus one.
The certified libraries have radii zero or one whenever repair succeeds, and
strict energy descent would give \(\rho\leq q\).  No structural or sharp
worst-case radius theorem is known.

\paragraph{Structural shape criteria.}
The active-carrier theorem is algorithmic.  The finite C32/C46/P49 panel does
not yield forbidden-pattern or shape-theoretic characterizations of support,
purity, or repair.

\paragraph{Enumeration.}
Exact participation types give FPT existential decision and compressed
classification.  A Delay-FPT algorithm for literal named outputs is not an
implemented result.

\paragraph{Continuous geometry.}
The C34 route would require a separate transcription and a continuous or
polygonal theorem.  It is costly, uncertain, and absent from the dependency
graph.

\paragraph{Algebraic completion of P21.}
The exact Markov degree, a complete Markov basis, four block-internal bases,
normality or compatible-projection hypotheses, and a structured lift remain
open.  They are deliberately parked: none is required for
\cref{thm:fpt,prop:p21-one,prop:p21-markov}.

\paragraph{Priority boundary.}
The primary-source audit located no parameter-preserving collision
with the precise active-carrier theorem, but it was not a systematic
MathSciNet/zbMATH and citation-forest search.  Any novelty, firstness, or
priority claim would require reopening that audit.  The current wording makes
no such claim.

\subsection{Conclusion}

Aggregate target margins contain enough information to define a useful
relaxation, but not enough to collapse exact coverage, repair, and mere
existence into one notion.  Largest-piece residual area restores algorithmic control:
after a largest pose is fixed, all residual activity lives in a
parameter-bounded carrier.  The finite libraries make the logical
distinctions concrete, while P21 shows that even complete knowledge of one
fiber does not substitute for relation-lattice or all-right-hand-side Markov
analysis.

\appendix
\section{Evidence and responsibility map}\label{app:evidence}

\subsection{Finite libraries}

\begin{center}
\small
\begin{tabular}{@{}L{0.20\textwidth}L{0.72\textwidth}@{}}
\toprule
Statement & Canonical repository evidence\\
\midrule
C32 source and normalization &
\path{registry/c32_source_lock.json};
\path{reports/P42_C32_SOURCE_AUDIT.md}.\\
C32 complete census &
\path{reports/P42_C32_EXACT_CENSUS.md};
Package B certificate, replay, and mutation outputs.\\
C46 source and normalization &
\path{registry/c46_source_lock.json};
\path{reports/P42_C46_SOURCE_AUDIT.md}.\\
C46 complete census &
\path{packages/package_a_c46_reproduction/realization/REPORT.md}
and its certificate, verifier, and mutation results.\\
P49 source specimen &
\path{packages/package_c_comparative_hypergraphs/specimens/SPECIMENS.md}.\\
Three-library fibers &
\path{reports/P42_COAREA_FPT_CLOSURE.md};
\path{results/coarea_fpt_analysis.json}.\\
\bottomrule
\end{tabular}
\end{center}

The source records establish displayed geometry and arrangements.  The shipped
artifacts establish discrete normalization, exhaustive realization counts,
margin fibers, graph components, and radii.  Raw copyrighted raster files
are intentionally absent; source-image audits require local bytes matching
the recorded SHA-256 digests.

\subsection{P21}

\begin{center}
\small
\begin{tabular}{@{}L{0.25\textwidth}L{0.67\textwidth}@{}}
\toprule
Statement & Canonical repository evidence\\
\midrule
Source identity and geometry &
\path{packages/package_c_comparative_hypergraphs/specimens/p21_locked.json}
and \path{packages/package_c_comparative_hypergraphs/specimens/SPECIMENS.md}.\\
72-variable quotient and repairs &
\path{reports/P42_ROBUST_P21_QUOTIENT_CLOSURE.md} and analysis JSON.\\
Support-three filtration &
\path{reports/P42_DEGREE3_AUGMENTATION_CLOSURE.md} and analysis JSON.\\
Support-four ceiling &
\path{reports/P42_SUPPORT4_CEILING_CLOSURE.md} and analysis JSON.\\
Minimum three-orbit generation &
\path{reports/P42_MINIMUM_ORBIT_GENERATORS_CERTIFICATE.json}.\\
Markov no-go and degree lower bound &
\path{reports/P42_MARKOV_NO_GO_CERTIFICATE.json};
\path{reports/P42_SINGLE_BLOCK_MARKOV_DEGREES.json}.\\
Independent integrated replay &
\path{reports/P42_ORBIT_MARKOV_CLOSURE_AUDIT.md}.\\
\bottomrule
\end{tabular}
\end{center}

\subsection{Claim boundary}

This paper asserts only the theorem and certified finite statements listed in
\path{docs/CLAIM_LEDGER.md}, at their stated scopes.  It makes no novelty,
firstness, or priority claim.  In particular, the exact P21 Markov degree and a
complete Markov basis remain open.

\section{Reproducibility contract}\label{app:reproduction}

All commands below are run from the repository root with Python in UTF-8,
no-bytecode mode and the exact environment in \path{requirements.txt}.

\subsection{Manifest and core replay}

\begin{verbatim}
python -X utf8 -B scripts/verify.py --profile manifest
python -X utf8 -B scripts/verify.py --profile core
\end{verbatim}

The core profile checks the complete frozen manifest, manuscript/PDF receipt,
C32 and C46 realization certificates, the P21/P26/P49 specimen locks, both
residual-area producers, an integrated P21 smoke replay, archive-level
checksums, import closure, and repository policy.

\subsection{Full mutation and P21 replay}

\begin{verbatim}
python -X utf8 -B scripts/verify.py --profile full
\end{verbatim}

The full profile adds all three mutation suites.  Its integrated P21 verifier
reconstructs the 72 participating columns, 7,444
support-four relations, 1,119 target orbits, rank-nine quotient, minimum
triples, and Markov obstructions without importing the principal
support-four analysis file.  It then repeats the full replay in a temporary
tree containing only the manifest allowlist; this detects hidden dependencies
on virtual environments, caches, build debris, or unshipped files.

\subsection{Manuscript verification}

The manuscript static verifier checks the frozen section order, labels,
citations, load-bearing theorem fragments, no-firstness boundary, and
presence of the built PDF.  The PDF is then reopened, text-extracted for
sanity, rendered to PNG, and visually inspected page by page.  Build and
verification commands are in \texttt{paper/BUILD.md}.

The code and result filenames beginning with \texttt{coarea\_} are stable
legacy artifact identifiers.  In the mathematical prose the parameter is
called the largest-piece residual area.

\begin{thebibliography}{99}

\bibitem{HermanKuba1999}
G.~T.~Herman and A.~Kuba (eds.),
\emph{Discrete Tomography: Foundations, Algorithms, and Applications},
Birkh\"auser, 1999;
doi:10.1007/978-1-4612-1568-4.

\bibitem{Sallows2011}
L.~Sallows,
\emph{Geometric magic squares},
Math. Intelligencer 33 (2011), 25--31;
doi:10.1007/s00283-011-9229-0.

\bibitem{SallowsGallery}
L.~Sallows,
\emph{Geomagic Squares: Gallery},
\url{https://www.geomagicsquares.com/gallery.php},
accessed 26 July 2026.

\bibitem{Cameron2011}
P.~J.~Cameron,
\emph{Geomagic squares},
Cameron Counts, 21 January 2011;
\url{https://cameroncounts.wordpress.com/2011/01/21/geomagic-squares/}.

\bibitem{ChrobakEtAl2003}
M.~Chrobak, C.~Couperus, C.~D\"urr, and G.~J.~Woeginger,
\emph{On tiling under tomographic constraints},
Theoret. Comput. Sci. 290 (2003), 2125--2136;
arXiv:cs/0108010.

\bibitem{ChrobakDurr2001}
M.~Chrobak and C.~D\"urr,
\emph{Reconstructing polyatomic structures from discrete X-rays:
NP-completeness proof for three atoms},
Theoret. Comput. Sci. 259 (2001), 81--98;
doi:10.1016/S0304-3975(99)00325-4.

\bibitem{DurrEtAl2000}
C.~D\"urr, E.~Goles, I.~Rapaport, and E.~R\'emila,
\emph{Tiling with bars under tomographic constraints},
Theoret. Comput. Sci. 290 (2003), 1317--1329;
arXiv:cs/9903020.

\bibitem{ChrobakEtAl2012}
M.~Chrobak, C.~D\"urr, F.~Gu\'i\~nez, A.~Lozano, and N.~K.~Thang,
\emph{Tile-packing tomography is NP-hard},
Algorithmica 64 (2012), 267--278;
doi:10.1007/s00453-011-9498-1.

\bibitem{FrosiniSimi2004}
A.~Frosini and P.~Simi,
\emph{The NP-completeness of a tomographical problem on bicolored domino
tilings},
Theoret. Comput. Sci. 319 (2004), 447--454;
doi:10.1016/j.tcs.2004.02.004.

\bibitem{FrosiniSimi2005}
A.~Frosini and P.~Simi,
\emph{The reconstruction of a subclass of domino tilings from two
projections},
Discrete Appl. Math. 151 (2005), 154--168;
doi:10.1016/j.dam.2005.02.032.

\bibitem{GritzmannLangfeld2020}
P.~Gritzmann and B.~Langfeld,
\emph{On polyatomic tomography over Abelian groups},
Discrete Comput. Geom. 64 (2020), 290--303;
doi:10.1007/s00454-020-00180-5.

\bibitem{FellowsEtAl2023}
M.~R.~Fellows, M.~Grobler, N.~Megow, A.~E.~Mouawad,
V.~Ramamoorthi, F.~A.~Rosamond, D.~Schmand, and S.~Siebertz,
\emph{On solution discovery via reconfiguration},
ECAI 2023, 700--707;
doi:10.3233/FAIA230334; arXiv:2304.14295.

\bibitem{TuckerFoltz2023}
J.~Tucker-Foltz,
\emph{Locked polyomino tilings},
arXiv:2307.15996.

\bibitem{GrossYamzon2021}
E.~Gross and N.~Yamzon,
\emph{Binomial ideals of domino tilings},
Discrete Math. 344 (2021), 112530;
arXiv:2008.02896.

\bibitem{Chin2019}
T.~L.~Chin,
\emph{A computational commutative algebra approach to tilings},
undergraduate honors thesis, Brown University, 2019;
\url{https://repository.library.brown.edu/studio/item/bdr:919199/}.

\bibitem{GuruswamiLin2019}
V.~Guruswami and P.~Lin,
\emph{Parameterized inapproximability of Exact Cover and Nearest Codeword},
arXiv:1905.06503.

\bibitem{DiaconisSturmfels1998}
P.~Diaconis and B.~Sturmfels,
\emph{Algebraic algorithms for sampling from conditional distributions},
Ann. Statist. 26 (1998), 363--397;
doi:10.1214/aos/1030563990.

\bibitem{AokiTakemuraYoshida2008}
S.~Aoki, A.~Takemura, and R.~Yoshida,
\emph{Indispensable monomials of toric ideals and Markov bases},
J. Symbolic Comput. 43 (2008), 490--507;
arXiv:math/0511290.

\bibitem{CharalambousThomaVladoiu2016}
H.~Charalambous, A.~Thoma, and M.~Vladoiu,
\emph{Binomial fibers and indispensable binomials},
J. Symbolic Comput. 74 (2016), 578--591;
arXiv:1501.05142.

\bibitem{RauhSullivant2019}
J.~Rauh and S.~Sullivant,
\emph{Lifting Markov bases and higher codimension toric fiber products},
J. Symbolic Comput. 74 (2016), 276--307;
arXiv:1404.6392.

\bibitem{CreignouEtAl2017}
N.~Creignou, A.~Meier, J.-S.~M\"uller, J.~Schmidt, and H.~Vollmer,
\emph{Parameterized enumeration with ordering},
Algorithmica 77 (2017), 202--223;
arXiv:1309.5009.

\end{thebibliography}


\end{document}
